\documentclass[10pt]{article}
\usepackage[margin=0.73in]{geometry}
\usepackage{amsmath,amssymb,amsthm,mathtools}
\usepackage{graphicx}
\usepackage{microtype}
\usepackage{xcolor}
\usepackage[colorlinks=true,linkcolor=blue!50!black,citecolor=blue!50!black,urlcolor=blue!50!black]{hyperref}

\newtheorem{theorem}{Theorem}
\newcommand{\DX}{\Delta_{\mathbb X}}
\newcommand{\Dp}{\mathbb D^+}
\newcommand{\Dm}{\mathbb D^-}
\newcommand{\ip}[2]{\left\langle #1,#2\right\rangle}
\setlength{\parindent}{0pt}
\setlength{\parskip}{4pt}

\title{The conjectured grid-gradient monotonicity can fail}
\author{Candidate counterexample for arXiv:2101.08108}
\date{27 August 2026}

\begin{document}
\maketitle

\vspace{-1.2em}
\begin{center}
\fbox{\parbox{0.94\textwidth}{\textbf{Status.}
Candidate counterexample, likely valid, pending expert review.  This concerns
only the fourth bullet of source Remark~4.2.}}
\end{center}

\section*{Source conjecture}

Bottazzi studies the hyperfinite Neumann grid flow
\begin{equation}\label{eq:flow}
 u_t=\DX\,{}^*\!\phi(u)
\end{equation}
and states four conjectural regularity properties in Remark~4.2.  The final
one asserts that the $L^2$ norm of
$\mathbb D^\pm{}^*\!\phi(u(t,\cdot))$ is nonincreasing in time
\cite[p.~18]{Bottazzi}.

\begin{center}
\includegraphics[width=0.92\textwidth]{figures/open_problem_crop.png}
\end{center}
\vspace{-0.6em}
\textbf{Transcription of the target bullet.} ``The $L^2$ norm of
$\mathbb D^\pm{}^*\!\phi(u(t,\cdot))$ is nonincreasing in time.''

\section*{Counterexample}

Define the standard function
\begin{equation}\label{eq:phi}
 \phi(s)=
 \begin{cases}
  4s,&s\leq0,\\
  s(2-s)^2,&0\leq s\leq2,\\
  0,&s\geq2.
 \end{cases}
\end{equation}
The values and first derivatives match at $0$ and $2$, so $\phi\in C^1(\mathbb
R)$ and is globally Lipschitz.  It satisfies $\phi(0)=0$ and $\phi(s)\geq0$
for $s\geq0$, as required in the source.  On the middle interval,
\begin{equation}\label{eq:dphi}
 \phi'(s)=(2-s)(2-3s),
\end{equation}
which is strictly negative for $2/3<s<2$ and vanishes at $s=2$.

Take $\Omega=(0,1)$ and its source hyperfinite grid.  Choose a nonzero
$C_c^\infty((1/4,3/4))$ bump, restrict its nonstandard extension to the
grid, and call the resulting internal grid function $\rho$.  After rescaling,
\[
 0\leq\rho\leq1,
 \qquad \operatorname{supp}\rho\subset[1/4,3/4],
\]
and set
\begin{equation}\label{eq:u0}
 u_0=2-\frac12\rho,\qquad v_0={}^*\!\phi(u_0).
\end{equation}
This is a nonnegative bounded initial datum satisfying source hypothesis
(5').  It is already a grid function, so the initial projection fixes it.
Its discrete derivatives have finite norms, and $v_0$ vanishes in a boundary
layer.

\begin{theorem}\label{thm:counterexample}
For the source solution of \eqref{eq:flow} with \eqref{eq:phi} and
\eqref{eq:u0}, both functions
\[
 t\longmapsto\|\Dp{}^*\!\phi(u(t,\cdot))\|_2,
 \qquad
 t\longmapsto\|\Dm{}^*\!\phi(u(t,\cdot))\|_2
\]
have positive right derivative at $t=0$.  In particular, the fourth
conjectural bullet of Remark~4.2 in \cite{Bottazzi} is false.
\end{theorem}

\section*{Proof Intuition}

The expected Lyapunov calculation contains an extra factor that determines
the sign.  If $v=\phi(u)$, grid integration by parts gives
\[
 \frac{d}{dt}\frac12\|\mathbb Dv\|_2^2
 =-\ip{\phi'(u)\DX v}{\DX v}.
\]
This is nonpositive on an increasing branch of $\phi$, but it is
\emph{positive} on a decreasing branch.  We place all nonconstant activity
inside such a branch and make $v$ zero near the boundary, eliminating every
possible boundary-convention issue.

\section*{Proof}

Put $v(t)={}^*\!\phi(u(t))$ and
\[
 E_+(t)=\frac12\|\Dp v(t)\|_2^2.
\]
Both $v_0$ and
$v_t(0)={}^*\!\phi'(u_0)u_t(0)$ vanish in a grid boundary layer: outside
$\operatorname{supp}\rho$ we have $u_0=2$, hence $v_0=0$ and
$\phi'(u_0)=0$.  Thus the source Neumann Laplacian agrees on the relevant
support with $\Dm\Dp$, and hyperfinite summation by parts has no boundary
term.  Differentiating exactly and using \eqref{eq:flow},
\begin{align}
 E_+'(0)
 &=\ip{\Dp v_0}{\Dp v_t(0)}
   =-\ip{\DX v_0}{v_t(0)} \notag\\
 &=-\ip{\DX v_0}{{}^*\!\phi'(u_0)\DX v_0}
   =-\int_{\Omega_{\mathbb X}}
       {}^*\!\phi'(u_0(x))|\DX v_0(x)|^2\,d\mathbb X.
 \label{eq:energy}
\end{align}

Where $\rho(x)>0$, equation \eqref{eq:u0} gives
$3/2\leq u_0(x)<2$, so \eqref{eq:dphi} gives
$\phi'(u_0(x))<0$.  Where $\rho(x)=0$, one has $u_0(x)=2$ and
$\phi'(u_0(x))=0$.  Therefore every summand in the final expression in
\eqref{eq:energy} is nonnegative.

At least one summand is strictly positive.  Indeed, $v_0$ is nonzero and
compactly supported in the connected grid.  If $\DX v_0$ vanished at every
point where $v_0>0$, applying the discrete maximum principle at a positive
maximum and propagating equality along its positive component would force
that component to meet a zero point with the same positive value, a
contradiction.  Hence
\[
 E_+'(0)>0.
\]
The initial gradient is nonzero, so its norm itself has positive right
derivative.  Replacing $\Dp$ by $\Dm$ gives the same energy by shifting the
hyperfinite sum; boundary terms again vanish.  This proves the theorem. \qed

\section*{Exact five-node check}

The sign is already visible with unit mesh on a five-node Neumann path.  Take
\[
 u=(2,2,3/2,2,2),\qquad
 v=\phi(u)=(0,0,3/8,0,0).
\]
Then
\[
 \Delta v=(0,3/8,-3/4,3/8,0),\qquad
 \frac12\|D^+v\|_2^2=\frac9{64},
\]
and direct differentiation gives
\[
 E_+'(0)=-\sum_j\phi'(u_j)(\Delta v_j)^2=\frac{45}{64}>0.
\]
The supplied verifier performs this calculation with exact rational
arithmetic.  Mesh scaling only multiplies the derivative by a positive power.

\section*{Verification, scope, and novelty}

The hypotheses on $\phi$ and $u_0$, the chain rule, the Neumann localization,
both difference directions, and the strict sign were audited separately.
This counterexample affects only the fourth bullet of Remark~4.2; the three
Dirac-Young-measure bullets and the paper's main theorems are untouched.

A bounded search on 27 August 2026 checked the run indexes, official arXiv
record, published Springer version, exact phrases from the conjecture, and
erratum/correction queries.  The published version retains the statement, and
no explicit correction or counterexample was found.  Novelty is plausible but
remains provisional pending a broader expert literature review.

\begin{thebibliography}{9}
\bibitem{Bottazzi}
E. Bottazzi,
\emph{Describing limits of integrable functions as grid functions of
nonstandard analysis}, Partial Differential Equations and Applications
\textbf{2} (2021), article 51; arXiv:2101.08108.
\end{thebibliography}

\end{document}
